\documentclass[10pt,a4paper]{amsart}
\usepackage[margin=19mm]{geometry}
\usepackage{amsmath,amssymb,mathtools}
\usepackage[scr=rsfs,cal=euler]{mathalfa}
\usepackage{xfrac}
\usepackage[unicode,hidelinks,bookmarks=false]{hyperref}
\newcommand{\ie}{i.e.\ }
\newcommand{\cf}{cf.\ }
\newcommand{\resp}{respectively\ }
\newcommand{\Subsection}{\subsection}
\providecommand{\nobreakdash}{\nobreak\hbox{-}}
\setlength{\emergencystretch}{2em}
\multlinegap0pt
\usepackage{bm}
\usepackage{bbm}
\usepackage{dsfont}
\usepackage{xspace}
\usepackage{cancel}
\usepackage{mathtools}
\usepackage{amsthm}
\makeatletter
\@ifundefined{theorem}{\newtheorem{theorem}{Theorem}}{}
\makeatother
\title[On graphs embeddable in a layer of a hypercube and their ext]{On graphs embeddable in a layer of a hypercube and their extremal numbers}
\author{Maria Axenovich, Ryan R. Martin,, Christian Winter}
\date{Source: arXiv:2303.15529v4 (CC BY 4.0)}
\begin{document}
\maketitle

\noindent\emph{Source and licence.} On graphs embeddable in a layer of a hypercube and their extremal numbers. Version arXiv:2303.15529v4 (CC BY 4.0), distributed under the Creative Commons Attribution 4.0 International licence, as recorded in the arXiv licence metadata for this exact version and shown on the abstract page https://arxiv.org/abs/2303.15529v4. Full rights notice: licenses/arxiv-2303.15529-CC-BY-4.0.txt.\par\smallskip

\noindent\textit{Editorial scope.} Counted unit: the Theorem whose source label is \texttt{None} in the article (source0.tex lines 1359--1363, source label \texttt{None}), delivered together with the article's own complete original proof of it (source0.tex lines 1364--1435, one \texttt{proof} environment of 72 lines). No mathematics is written, repaired or re-derived: statement and proof are the article's own text line for line. Required context retained uncounted: none. Every definition, display and label of those lines is retained unchanged. Omitted from this package and not redistributed: 1--1358, 1436--1472. Nothing inside a retained excerpt is omitted or altered: every retained body is delivered exactly as the article's lines are written, so no omission receipt of any kind accompanies this package. Citations: the retained excerpts carry no citation command at all, so no bibliography entry of the article is transcribed and no reference is redistributed. Reference adaptation: none was needed and none was made; every reference inside the delivered unit resolves inside it, so no unresolved reference or citation is produced. Printed slips kept literal: no printed word, symbol, hypothesis or number of the counted statement or of its proof was altered. Layout and class: the article's own class is \texttt{article} with packages including amsmath, amsthm, amsfonts, amssymb, setspace, fullpage, enumitem, graphicx, floatpag, xcolor, amsrefs, float, caption, subcaption; the wrapper is the lane's pinned \texttt{amsart} editorial wrapper at 10pt on A4, which restates the theorem-like environments the retained text needs and adds the article's own macro definitions that the retained text needs (none), with extra wrapper packages (bm, bbm, dsfont, xspace, cancel, mathtools). The delivered numbering is the unit's own. Assets: no retained span contains a figure environment, a graphics-inclusion command, a PDF-page inclusion, a table or a TikZ picture; no image file is copied into this package, the package declares no asset dependency and no image file is redistributed with it. Licence: version arXiv:2303.15529v4 of this article is distributed under the Creative Commons Attribution 4.0 International licence, as recorded in the arXiv licence metadata for this exact version and shown on the abstract page https://arxiv.org/abs/2303.15529v4. A full rights notice is delivered at licenses/arxiv-2303.15529-CC-BY-4.0.txt. Characters: every retained excerpt is pure ASCII, so the delivered engine, XeLaTeX with its default Latin Modern text font, reports no missing character.
\begin{theorem}
For any integer $m\geq 2$ and any positive integer $n$, there exist an integer $N$ and  a function $F: V(Q_n) \rightarrow V(Q_N)$, such that 
for any two vertices $u$ and $v$ which are adjacent in $Q_n$,  $d_H(F(u), F(v))=m$ and $F$ maps all vertices of $Q_n$ either in one vertex layer of $Q_N$ (if $m$ is even) or 
in two consecutive vertex layers of $Q_N$ (if $m$ is odd).
\end{theorem}

\begin{proof}

Here, we shall present functions $f, f', f_k$ mapping $V(Q_n)$ into the vertex set of some larger hypercube, for a fixed $k\in\mathbb{N}\cup \{0\}$  such that 
for any two adjacent in $Q_n$ vertices $u$ and $v$,  $d_H(f(u), f(v))= d_H(f'(u), f'(v))=3$ and $d_H(f_k(u), f_k(v))=2k+2$. 
Moreover, both $f$ and $f'$ map vertices of $Q_n$ into two consecutive vertex layers  and $f_k$ maps vertices of $Q_n$ into one layer. We shall then define $F$ based on one of the functions $f, f'$, or $f_k$.\\

For any vector $w$, let  $w[i]$ denote the $i$th component of $w$ and $||w||$ denote the number of $1$'s in $w$. \\


Let $[(2k+2)n]$ be split into $n$ consecutive intervals of length $2k+2$.   For a binary vector $w$ of length $(2k+2)n$, let $w[[i]]$ be $w$ restricted to the $i$th interval of length $2k+2$. Formally, $w[[i]] = w[(2k+2)(i-1)+1] w[(2k+2)(i-1)+2] \cdots w[(2k+2)i]$. We define $f_k: V(Q_n) \rightarrow V(Q_{(2k+2)n})$ as follows:

\begin{equation} \nonumber
f_k(v)[[i]] = \begin{cases}
0101\cdots 01,  \mbox{ if }  v[i]=0, \\
1010\cdots 10,  \mbox{ if }  v[i]=1.
\end{cases}
\end{equation}



\vskip 1cm


Let $[2n+1]$ be split into $n$ consecutive intervals of length $2$ and one last element. For a binary vector $w$ of length $2n+1$, let $w[[i]]$ be a triple corresponding to $w$ restricted to the $i$th interval and the last element, i.e., $w[[i]]= w[2i-1] w[2i] w[N']$.   We define $f: V(Q_n)\rightarrow V(Q_{2n+1})$ as follows:

\begin{equation}\nonumber
f(v)[[i]] = \begin{cases}
010,  \mbox{ if }  v[i]=0  \mbox{ and  }  ||v|| \mbox{ is even, }  \\
100,  \mbox{ if }  v[i]=1  \mbox{ and  }  ||v|| \mbox{ is even, }  \\
011,  \mbox{ if }  v[i]=0  \mbox{ and  }  ||v|| \mbox{ is odd, }  \\
101,  \mbox{ if }  v[i]=1  \mbox{ and  }  ||v|| \mbox{ is odd. }  \\
\end{cases}
\end{equation}

It is clear here that if $u$ and $v$ are adjacent in $Q_n$, the images $f(u)$ and $f(v)$ are at Hamming distance $3$.
		
	
\vskip 1cm	

Let $[3n]$ be split into $n$ consecutive intervals of length $3$.  For a binary vector $w$ of length $3n$, let $w[[i]]$ be a triple corresponding to $w$ restricted to the $i$th interval  $w[[i]]= w[3i-2] w[3i-1] w[3i]$.  We define $f':V(Q_n)\rightarrow V(Q_{3n})$ as follows:


Let 
\begin{equation} \nonumber
f'(v)[[i]] = \begin{cases}
010,  \mbox{ if }  v[i]=0, \\
100,  \mbox{ if }  v[i]=1 ~ \mbox{ and   } ||v|| \mbox{ is even}, \\
101,  \mbox{ if }  v[i]=1,  v[j]=0 \mbox{ for any } j>i, ~ \mbox{ and   } ||v|| \mbox{ is odd},\\
100,  \mbox{ if }  v[i]=1,  v[j]=1 \mbox{ for some } j>i, ~ \mbox{ and   } ||v|| \mbox{ is odd}.
\end{cases}
\end{equation}		
		
Assume that $v$ and $u$ are adjacent in $Q_n$ and differ in position $i$ such that $v$ is zero in this position.
We shall verify that the distance between $f(v)$ and $f(u)$ is $3$.
Note that $f(v)$ and $f(u)$ coincide in all triples corresponding to $0$'s of $u$. Moreover, they coincide on those triples $\ell$, where $u[\ell]=1$, $j\neq i$, and $\ell$ is not a position of the last $1$ of $u$ or $v$.  Let $j$ be the last position of $1$ in $u$. Note that $i$ could be equal to $j$.  \\

If $w(v)$ is even, then $w(u)$ is odd and $f(u)[[j]]=101$. If $i=j$, then $f(v)[[j]]=010$ and $f(u)$ and $f(v)$ coincide in all other triples.  If $i<j$,  then $u[j]=v[j] =1$, $f(u)[[j]]=101$, $f(v)[[j]]=100$, 
$f(u)[[i]]=100$, and $f(v)[[i]]=010$. On all other triples $f(v)$ and $f(u)$ coincide. We see that in both cases $f(u)$ and $f(v)$ are at distance $3$.\\

If $w(v)$ is odd, then $w(u)$ is even and $f(u)[[i]]=100$. 
Let $k$ be the last position of $1$ in $v$. So, $f(v)[[k]]= 101$. 
We also have $f(u)[[k]]=100$ and $f(v)[[i]]= 010$. Then $f(u)$ and $f(v)$ are at distance $3$.\\~\\


Now, let $m\geq 2$ be given. If $m$ is even, let $m=2k+2$, for non-negative integer $k$. Then let $F(u) = f_k(u)$ for any $u\in V(Q_n)$. 
If $m$ is odd and $m=3$ let $F(u) = f(u)$ for any $u\in V(Q_n)$. 
If $m$ is odd and $m= 3 +2\ell$ for some positive integer $\ell$, we define $F$ by considering either $f$ or $f'$ and adjusting $2\ell$ coordinates to each embedded vertex 
that are $0\cdots 0 1 \cdots 1$ or $ 1\cdots 1 0\cdots 0$, depending whether the vertex is embedded in one layer of the other. 
Formally, in case of $f$, for example, let  $N= 2n+1+ 2\ell$ and let for any vertex $u$ of $Q_n$, $F(u)$ restricted to the first $2n+1$ coordinates be $f(u)$. 
In addition, if $w(u)$ is even, let the last $2\ell$ coordinates of $F(u)$ be $0\cdots 0 1 \cdots 1$ and 
if $w(u)$ is odd, let the last $2\ell$ coordinates of $F(u)$ be $1\cdots 1 0 \cdots 0$,  with $\ell$ $0$'s and $\ell$ $1$'s respectively. 
\end{proof}

\end{document}
